
\section*{7. Appendice : Spécialisation en théorie des intersections}
\begin{center}
\underline{par A. Grothendieck}
\end{center}

\textbf{Lemme 7.1.} \uline{Soit} \(i:Y\to X\) \uline{une immersion fermée régulière} (VII 1.4)\uline{, de codimension bornée, de sorte que} \(i\) \uline{est un morphisme parfait} (VII 1.9)\uline{, et qu'on peut donc définir le morphisme} (IV 2.12)
\[
 i^*:K_\bullet(X)\longrightarrow K_\bullet(Y).
\]
\uline{Soit d'autre part} \(\mathbb N\) \uline{le faisceau conormal pour} \(i\)\uline{, qui est un faisceau localement libre de type fini sur} \(Y\) (VII 1.6)\uline{, de sorte qu'on peut former}
\[
 \lambda_{-1}(\mathbb N)\in K^{\bullet}(Y).
\]
\uline{Considérons aussi l'nomomorphisme image directe} (IV 2.11)
\[
 i_*:K_\bullet(Y)\longrightarrow K_\bullet(X).
\]
\uline{Ceci posé, on a la formule}
\[
 i^*i_*(x)=x\lambda_{-1}(\mathbb N)\qquad\text{pour } x\in K_\bullet(Y),
\]
\uline{où le produit du deuxième membre est relatif à la structure de} \(K^{\bullet}(Y)\)\uline{-module de} \(K_\bullet(Y)\) \uline{définie dans} (IV 2.10).

La démonstration est, essentiellement sans changement, la même que pour la même formule, écrite en travaillant avec \(K^{\bullet}(X)\), \(K^{\bullet}(Y)\), établie dans (VII 2.7).

\textbf{Corollaire 7.2.} \uline{Soient} \(f:X\to S\) \uline{un morphisme plat, avec} \(S\) \uline{un trait i.e. le spectre d'un anneau de valuation discrète} \(A\), \(s\) \uline{le point fermé de} \(S\), \(t\) \uline{son point générique,} \(i:X_s\to X\) \uline{l'inclusion canonique. Alors} \(i\) \uline{est de tor-dimension finie donc} (IV 2.12) \uline{on peut définir}
\[
 i^*:K_\bullet(X)\longrightarrow K_\bullet(X_s),
\]
